%===========================================================================
\chapter{RISC-V Pipelined Processor Integration}\label{ch:cpu}
%===========================================================================

A MAC unit is only useful inside a processor that fetches, decodes and
issues its instructions. To show this, the MAC unit was built into a complete
\textbf{5-stage pipelined RISC-V processor} (\code{rtl/rv32\_core.v}, 638
lines) that implements RV32I, the multiply subset of RV32M, and the eight MAC
instructions of Chapter~\ref{ch:spec}. The processor runs real RISC-V machine
code produced by a purpose-written assembler. Its results are checked
bit-exactly in simulation (Chapter~\ref{ch:sysres}) and it was synthesised and
placed-and-routed on an FPGA (Chapter~\ref{ch:synth}).

\section{Pipeline organisation}
Figure~\ref{fig:cpu} shows the datapath. It is the classic five-stage RISC
pipeline (IF, ID, EX, MEM, WB) described by Patterson and
Hennessy~\cite{patterson}, extended with the MAC unit.

\begin{figure}[H]
  \centering
  \resizebox{\linewidth}{!}{% ===========================================================================
% fig_pipeline.tex - 5-stage RV32IM pipeline with the integrated MAC unit
% (included from report.tex)
% ===========================================================================
\begin{tikzpicture}[
    font=\scriptsize, >=Stealth, line width=0.5pt,
    blk/.style={draw, rounded corners=1.5pt, fill=accent!10, align=center, minimum height=0.8cm},
    mac/.style={draw, rounded corners=1.5pt, fill=orange!25, align=center},
    preg/.style={draw, fill=gray!30, minimum width=0.32cm, minimum height=6.2cm},
    unit/.style={draw, dashed, rounded corners=1.5pt, fill=green!8, align=center},
    lab/.style={font=\bfseries\footnotesize, text=accent},
  ]
  % ---------------- pipeline registers ----------------
  \node[preg] (r1) at (2.55,0) {};
  \node[preg] (r2) at (6.25,0) {};
  \node[preg] (r3) at (10.35,0) {};
  \node[preg] (r4) at (13.25,0) {};
  \node[font=\tiny, above=1pt of r1] {IF/ID};
  \node[font=\tiny, above=1pt of r2] {ID/EX};
  \node[font=\tiny, above=1pt of r3] {EX/MEM};
  \node[font=\tiny, above=1pt of r4] {MEM/WB};

  % ---------------- stage labels ----------------
  \node[lab] at (1.1,4.05)  {IF};
  \node[lab] at (4.4,4.05)  {ID};
  \node[lab] at (8.3,4.05)  {EX};
  \node[lab] at (11.8,4.05) {MEM};
  \node[lab] at (14.2,4.05) {WB};

  % ---------------- IF ----------------
  \node[blk, minimum width=0.8cm] (pc) at (0.35,1.6) {PC};
  \node[blk, minimum width=1.3cm, minimum height=1.4cm] (imem) at (1.45,0.2) {Instr.\\memory};
  \node[blk, minimum width=0.7cm, minimum height=0.5cm] (add4) at (1.45,1.75) {+4};
  \draw[->] (pc.south) |- (imem.west);
  \draw[->] (pc.east) -- (add4.west);
  \draw[->] (imem.east) -- (imem.east -| r1.west);

  % ---------------- ID ----------------
  \node[blk, minimum width=1.4cm] (dec) at (3.6,1.9) {Decoder\\+ imm gen};
  \node[blk, minimum width=1.4cm, minimum height=1.6cm] (rf) at (3.6,-0.2) {Register\\file\\$32\times32$};
  \node[unit, minimum width=1.5cm] (haz) at (4.55,-2.35) {Hazard unit\\load/MAC-use\\stall};
  \draw[->] (r1.east |- dec.west) -- (dec.west);
  \draw[->] (r1.east |- rf.west) -- (rf.west);
  \draw[->] (dec.east) -- (dec.east -| r2.west);
  \draw[->] (rf.east) -- (rf.east -| r2.west);

  % ---------------- EX ----------------
  \node[blk, minimum width=0.55cm, minimum height=1.3cm] (fmux) at (7.0,0.6) {fwd\\mux};
  \node[blk, minimum width=1.6cm, minimum height=1.3cm] (alu) at (8.45,0.9) {ALU\\+ RV32M\\multiplier};
  \node[blk, minimum width=1.4cm] (br) at (8.45,2.55) {Branch\\compare};
  \node[mac, minimum width=1.6cm, minimum height=1.05cm] (mac1) at (8.45,-1.0) {MAC stage 1\\Booth $34\times34$};
  \node[unit, minimum width=1.5cm] (fwd) at (7.6,-2.6) {Forwarding\\unit};
  \draw[->] (r2.east |- fmux.west) -- (fmux.west);
  \draw[->] (fmux.east |- alu.west) -- (alu.west);
  \draw[->] (fmux.north) |- (br.west);
  \draw[->] (fmux.south) |- (mac1.west);
  \draw[->] (alu.east) -- (alu.east -| r3.west);
  \draw[->] (mac1.east) -- node[above, font=\tiny]{s1\_prod} (mac1.east -| r3.west);
  \draw[->, red!70!black] (br.north) -- ++(0,0.55) -| node[pos=0.25, above, font=\tiny]{taken: PC $\leftarrow$ target, flush IF/ID + ID/EX} (pc.north);

  % ---------------- MEM ----------------
  \node[blk, minimum width=1.5cm, minimum height=1.4cm] (dmem) at (11.8,0.9) {Data\\memory\\(byte/half/word)};
  \node[mac, minimum width=1.5cm, minimum height=1.05cm] (mac2) at (11.8,-1.0) {MAC stage 2\\ACC$_{64}$ + sat};
  \draw[->] (r3.east |- dmem.west) -- (dmem.west);
  \draw[->] (r3.east |- mac2.west) -- (mac2.west);
  \draw[->] (dmem.east) -- (dmem.east -| r4.west);
  \draw[->] (mac2.east) -- node[above, font=\tiny]{rd\_data} (mac2.east -| r4.west);

  % ---------------- WB ----------------
  \node[blk, minimum width=0.6cm, minimum height=1.6cm] (wbmux) at (14.2,0.0) {WB\\mux};
  \draw[->] (r4.east |- wbmux.west) -- (wbmux.west);
  % write-back to register file
  \draw[->, accent] (wbmux.south) -- ++(0,-2.95) -| node[pos=0.2, below, font=\tiny]{write-back (rd, value) \quad and WB $\rightarrow$ EX forwarding} (rf.south);

  % forwarding paths (dashed)
  \draw[->, dashed, green!45!black] (r3.south) ++(0,0.2) |- ($(fwd.east)+(0,0.12)$) node[pos=0.75, above, font=\tiny]{EX/MEM};
  \draw[->, dashed, green!45!black] (fwd.north) -- (fwd.north |- fmux.south);
  \draw[->, dashed, green!45!black] (haz.north) -- ++(0,0.5) -| node[pos=0.3, left, font=\tiny]{stall} (r1.south);

  % MAC span box
  \draw[orange!80!black, rounded corners, densely dotted, line width=0.8pt]
      ($(mac1.south west)+(-0.15,-0.15)$) rectangle ($(mac2.north east)+(0.15,0.15)$);
  \node[font=\tiny, text=orange!60!black, anchor=north west] at ($(mac1.south west)+(-0.1,-0.22)$)
      {mac\_unit: its two internal registers coincide with EX/MEM and MEM/WB};
\end{tikzpicture}
}
  \caption{5-stage RV32IM pipeline with the integrated MAC unit. Grey bars are
    pipeline registers; dashed green blocks are the hazard and forwarding
    units; the MAC unit (orange) spans EX and MEM.}
  \label{fig:cpu}
\end{figure}

\begin{table}[H]
  \centering
  \caption{Work done in each pipeline stage.}
  \label{tab:stages}
  \small
  \begin{tabularx}{\linewidth}{@{}lXX@{}}
    \toprule
    Stage & Ordinary instructions & MAC instructions \\
    \midrule
    IF  & PC register addresses the instruction memory; next PC = PC+4 or branch target & same \\
    ID  & decode, immediate generation, register-file read (with write-first bypass), hazard check & same; only \code{mac/msu/mul.acc/macu} read \code{rs1}/\code{rs2} \\
    EX  & forwarding muxes, ALU, RV32M multiplier, branch/jump resolution & forwarded operands enter \textbf{MAC stage 1}: decode + Booth multiply \\
    MEM & load/store access with byte/half/word alignment & \textbf{MAC stage 2}: accumulate, saturate, read-out \\
    WB  & write ALU result or load data to \code{rd} & write the MAC result to \code{rd} \\
    \bottomrule
  \end{tabularx}
\end{table}

\paragraph{Why the MAC unit fits without stalls.}
The MAC unit has two internal pipeline registers (Chapter~\ref{ch:impl}). In
the processor they line up exactly with the EX/MEM and MEM/WB registers: a MAC
instruction is in MAC stage~1 while it is in EX, in MAC stage~2 while it is in
MEM, and its \code{rd\_data\_o} is valid when it reaches WB. The MAC unit
never stalls, and the core never stalls EX, MEM or WB, so the two pipelines
always move together. The testbench checks this alignment on every MAC
write-back.

\section{Hazard handling}
\paragraph{Data hazards: forwarding.}
An instruction in EX can read a value that has been computed but not yet
written to the register file. The forwarding unit selects each EX operand
from three sources, in priority order:
\begin{equation}
  \text{op}_A =
  \begin{cases}
    \text{EX/MEM.result} & \text{if EX/MEM writes } rd = rs_1 \ne 0 \text{ and is not a load or MAC},\\
    \text{WB value}      & \text{else if MEM/WB writes } rd = rs_1 \ne 0,\\
    \text{ID/EX.rs1\_val} & \text{otherwise,}
  \end{cases}
\end{equation}
and the same for $\text{op}_B$. The register file is also bypassed in ID, so
an instruction three positions behind a writer always reads the new value.

\paragraph{Load-use and MAC-use stalls.}
A load result is ready only at the end of MEM, and so is a MAC result, because
it leaves MAC stage~2 at the MEM/WB boundary. An instruction that uses such a
result immediately afterwards must therefore wait one cycle. The hazard unit
detects
\begin{equation}
\begin{split}
  \text{stall} = {} & \text{IF/ID.valid} \wedge \text{ID/EX.valid} \wedge
                      (\text{ID/EX.load} \vee \text{ID/EX.mac}) \\
                    & \wedge \big(rd_{EX}=rs_{1,ID} \vee rd_{EX}=rs_{2,ID}\big) \wedge rd_{EX}\ne 0,
\end{split}
\end{equation}
freezes PC and IF/ID, and inserts a bubble into ID/EX. One cycle later, the
value is forwarded from WB. Section~\ref{sec:pipeviz} shows this happening in
the simulation.

\paragraph{Control hazards.}
Branches and jumps are resolved in EX. A taken branch loads the target into
the PC and turns the two younger instructions (in IF and ID) into bubbles, so
it costs two cycles. ECALL/EBREAK stops instruction fetch; when it reaches WB,
the core halts (\code{halted\_o}). An illegal instruction halts the core with
\code{error\_o}.

\section{Instruction set and system}
\begin{table}[H]
  \centering
  \caption{Instructions implemented by \code{rv32\_core}.}
  \label{tab:cpuisa}
  \small
  \begin{tabularx}{\linewidth}{@{}lX@{}}
    \toprule
    Group & Instructions \\
    \midrule
    RV32I arithmetic/logic & \code{ADD SUB SLL SLT SLTU XOR SRL SRA OR AND} and immediate forms \code{ADDI SLTI SLTIU XORI ORI ANDI SLLI SRLI SRAI}, \code{LUI AUIPC} \\
    RV32I control & \code{JAL JALR BEQ BNE BLT BGE BLTU BGEU} \\
    RV32I memory & \code{LB LH LW LBU LHU SB SH SW}, \code{FENCE} (no-op) \\
    RV32I system & \code{ECALL EBREAK} (halt) \\
    RV32M (multiply subset) & \code{MUL MULH MULHSU MULHU}, using one shared 33$\times$33 signed multiplier \\
    Xmac (custom-0) & \code{mac msu mul.acc macu rdacc.lo rdacc.hi rdacc.q15 clracc} \\
    \bottomrule
  \end{tabularx}
\end{table}

\noindent Division (\code{DIV}/\code{REM}) and CSR instructions are not
implemented, since none of the DSP programs need them. \code{riscv\_soc.v}
connects the core to a 4\,KiB instruction memory at address
\code{0x0000\_0000} and a 16\,KiB byte-enabled data memory at
\code{0x0001\_0000}. Both are initialised by the testbench with
\code{\$readmemh}.

\section{Software toolchain}
There is no RISC-V GCC toolchain on the development machine, and a standard
toolchain would not know the custom instructions anyway. So a small two-pass
assembler, \code{sw/rvasm.py} (386 lines of Python), was written for this
project. It supports all instructions in Table~\ref{tab:cpuisa}, the common
pseudo-instructions (\code{li la mv j call ret beqz bnez}, \ldots), labels,
\code{.equ}/\code{.include} and data directives. Its output encodings were
checked against reference encodings from the RISC-V specification (for
example \code{lw x5,8(x2)} $\rightarrow$ \code{0x00812283} and
\code{bne x1,x2,-8} $\rightarrow$ \code{0xfe209ce3}).
\code{sw/build.py} designs the test filter, computes the golden results and
assembles all the programs (Table~\ref{tab:progs}).

\begin{table}[H]
  \centering
  \caption{Programs executed on the processor.}
  \label{tab:progs}
  \small
  \begin{tabularx}{\linewidth}{@{}lrX@{}}
    \toprule
    Program & Static instr. & Purpose \\
    \midrule
    \code{isa\_test} & 282 & 18 self-checking tests: every instruction class, sign/zero extension, all branch types, call/return, RV32M, every MAC instruction, forwarding, load-use and MAC-use hazards, \code{x0} \\
    \code{fir\_mac\_loop} & 27 & 16-tap Q15 FIR over 128 samples using \code{mac} and \code{rdacc.q15} \\
    \code{fir\_rv32im\_loop} & 51 & the same filter in plain RV32IM (\code{mul}/\code{mulh} + 64-bit add with carry + software saturation) \\
    \code{fir\_mac\_unrolled} & 64 & MAC version with the tap loop fully unrolled \\
    \code{fir\_rv32im\_unrolled} & 166 & RV32IM version with the tap loop fully unrolled \\
    \bottomrule
  \end{tabularx}
\end{table}

Listing~\ref{lst:firloops} compares the two inner loops. Per tap, the MAC
version needs 7 instructions where RV32IM needs 12, because one \code{mac}
replaces the six-instruction 64-bit multiply-accumulate sequence. After the
loop, the MAC version converts and saturates the result with a single
\code{rdacc.q15}, where RV32IM needs a 14-instruction shift-and-saturate
routine. Both loops are scheduled so that no load-use stall occurs.

\begin{figure}[H]
\begin{minipage}[t]{0.40\linewidth}
\begin{lstlisting}[language={}, basicstyle=\ttfamily\scriptsize, numbers=none, title={\code{fir\_mac\_loop} (7 instr./tap)}]
inner:
  lh   t3, 0(t0)    # x
  lh   t4, 0(t1)    # h
  addi t0, t0, 2
  addi t1, t1, 2
  addi t2, t2, -1
  mac  x0, t3, t4   # ACC += x*h
  bnez t2, inner
  rdacc.q15 t5      # Q30->Q15, sat
\end{lstlisting}
\end{minipage}\hfill
\begin{minipage}[t]{0.56\linewidth}
\begin{lstlisting}[language={}, basicstyle=\ttfamily\scriptsize, numbers=none, title={\code{fir\_rv32im\_loop} (12 instr./tap)}]
inner:
  lh   t3, 0(t0)
  lh   t4, 0(t1)
  addi t0, t0, 2
  addi t1, t1, 2
  mul  a2, t3, t4    # product low
  mulh a3, t3, t4    # product high
  addi t2, t2, -1
  add  s4, s4, a2    # 64-bit add
  sltu a4, s4, a2    #   carry
  add  s5, s5, a3
  add  s5, s5, a4
  bnez t2, inner
  ... 14 instr.: acc>>>15 + saturate
\end{lstlisting}
\end{minipage}
\captionof{lstlisting}{FIR inner loops with and without the MAC unit.}
\label{lst:firloops}
\end{figure}

\section{DSP test case: low-pass FIR filter}
The benchmark is a realistic filtering task. A 16-tap low-pass filter
(windowed sinc, Hamming window, cut-off 0.1 cycles/sample, unity DC gain) is
quantised to Q15:
\[
h = \{-114, -159, -139, 291, 1450, 3284, 5246, 6524, 6524, 5246, 3284, 1450, 291, -139, -159, -114\}.
\]
The input contains 128 Q15 samples of a wanted low-frequency tone
($0.5\sin(2\pi\,0.015\,n)$) plus an unwanted high-frequency tone
($0.3\sin(2\pi\,0.3\,n)$). \code{build.py} computes the bit-exact expected output
$y[n]=\mathrm{sat}_{16}\big((\sum_k h[k]x[n-k]) \ggg 15\big)$ in Python
(independently of the hardware) and writes it to \code{build/fir\_golden.hex}
for the testbench.
